\section{Introduction}
Designing Internet of Things (IoT) networks often demands establishing connection between large number of devices, which can be achieved by activating communication interfaces at each of them. However, the devices present in the network may vary according to the type of interfaces they have at their availability and the amount of resources a given interface consumes. For example some device may be able to utilize Bluetooth and Wi-Fi connections, while the others may provide Zigbee and Z-Wave interfaces. Since IoT devices are typically resource-constrained, it may be desirable to try to establish every relevant connection, while minimizing the amount of resources required to do so. 

\begin{figure}[t]
\centering
\input{figures/combined-illustration.tex}
\caption{From left to right: a common input graph, a covering assignment, and a connecting assignment for the same instance.}
\label{fig:example_problems}
\end{figure}
We model the above problem as a graph problem. Every device in the network is represented by a vertex with a set of labels encoding its available interfaces. Two devices can communicate if they share a common available interface and are connected by an edge, representing that they are within communication range. A connection is established between two adjacent devices whenever they share a common active interface. 
We want to find an assignment of active interfaces to vertices so that we ensure either the \coverage or \connectivity property. 

In the \coverage problem, we wish to establish every possible direct connection in the network which means that both endpoints of every edge must share a common activated interface (see the middle panel in \Cref{fig:example_problems} for an example) while in the \connectivity problem, we require all vertices to be (not necessarily directly) connected in the subgraph induced by active interfaces (see the right panel in \Cref{fig:example_problems} for an example). In both problems, we assume that activating an interface induces a cost that depends on the type of interface and on the vertex that activates it. The goal is to minimize the max cost, i.~e., the maximum cost of  interfaces activated at a given vertex across all vertices. This cost criterion is meant to spread the energy consumption somewhat evenly between devices, ensuring the sustainability of the whole network.

\subsection{Related Work}

In the following, let $k$ denote the number of different types of interfaces across the network, $n$ the number of vertices in the input graph, $m$ the number of edges, $\Delta$ the maximum degree, $c_{\max}$ the maximum cost and $c_{\min}$ the minimum cost of an interface. In the \emph{homogeneous} setting, the costs only depend on the interface while in the \emph{heterogeneous} case they depend on the interface and on the vertex. 

\subparagraph*{Coverage Problem.}
The algorithmic study of problems in multi-interface networks started with the problem of covering a multi-interface network, introduced in Caporuscio et al. \cite{EnergeticPerformanceOfServiceOrientedMultiRadioNetworks}. 
This problem has been mainly studied in the homogeneous setting, both for min-sum \cite{CostMinimizationInWirelessNetworks} and min-max \cite{MinimizeTheMaximumDutyInMultiInterfaceNetworks} objective functions. There has been recent work on the heterogeneous case for the problem of maximizing some profit function while keeping the energy cost within budget \cite{OnBudgetConstraineCoverageInMultiInterfaceNetworksBranchwidthAndTreewidthPerspectives}.
The min-sum variant is APX-hard for any constant $k \geq 3$ and unit costs, and it is hard to approximate within an $o\br{\ln k}$ factor even with unit costs, but it has a $\min \brc{k-1, \sqrt{n} \br{1+\ln n}}$-approximation \cite{CostMinimizationInWirelessNetworks}.
The min-max variant is \NPhard for any constant $\Delta \geq 5$, $k \geq 16$ and for unit costs, and it is not approximable within a $\eta \ln \Delta$ factor for some constant $\eta$ even on trees and with unit costs. On the positive side, using an approximation of \textsc{Set Cover}, it admits a $\br{\ln \Delta +1 +b \cdot \min \brc{c_{\max}, \ln \Delta +1}}$-approximation, where $b$ is a parameter that depends on the structural properties of the input graph and is bounded by many other graph structutal parameters such as treewidth, arboricity or maximum degree. It also has a $(1+\br{k-2} \frac{c_{\max}}{2 c_{\min}})$-approximation and a $(\frac{\Delta}{2} \frac{c_{\max}}{c_{\min}})$-approximation \cite{MinimizeTheMaximumDutyInMultiInterfaceNetworks}. 

\subparagraph*{Connectivity Problem.} The problem of connecting a multi-interface network was introduced by Kosowski et al. \cite{ExploitingMultiInterfaceNetworksConnectivityAndCheapestPaths}. It was first studied in the homogeneous setting, both for the min-sum \cite{ExploitingMultiInterfaceNetworksConnectivityAndCheapestPaths} and the min-max objective functions \cite{MinimizeTheMaximumDutyInMultiInterfaceNetworks}.
The min-sum variant where all costs are homogeneous is APX-hard for $k \geq 2$, $\Delta \geq 4$ and unit costs already, but has a $2$-approximation \cite{ExploitingMultiInterfaceNetworksConnectivityAndCheapestPaths}. This result was improved by Athanassopoulos et al. \cite{EnergyEfficientCommunicationInMultiInterfaceWirelessNetworks} to a randomized $(\frac{3}{2} + \epsilon)$-approximation for any constant $\epsilon$, using a reduction to a \textsc{Minimum Spanning Tree} problem on hypergraphs. Moreover, this approximation algorithm applies to the more general problem where for each interface, the set of edges that it can cover is specified in the input.

The min-max variant where all costs are homogeneous is \NPhard for any constant $\Delta \geq 3$, $k \geq 10$ and for unit costs. Similarly to the coverage, it is not approximable within a $\eta \ln n$ factor even for unit costs, but has simple $(1+(k-2) \frac{c_{\max}}{2 c_{\min}})$- and $\frac{\Delta}{2} \frac{c_{\max}}{c_{\min}}$-approximations \cite{MinimizeTheMaximumDutyInMultiInterfaceNetworks}.

Athanassopoulos et al. \cite{EnergyEfficientCommunicationInMultiInterfaceWirelessNetworks} expand the study of the min-sum variant to the case where only certain vertices, called \textit{terminals} need to be spanned and there may be several terminal groups. They provide a randomized $4$-approximation for this case. They also initiate the study of this problem in the heterogeneous setting. They show that for heterogeneous costs, the min-sum variant is hard to approximate within $o(\ln n)$ even when all vertices are terminals, but the min-sum group \textsc{connectivity} heterogeneous case admits a $O(\ln n)$-approximation.

\subparagraph*{Other Problems.} There has been some work on the study of the coverage problem from the perspective of parameterized complexity, possibly with an additional constraint on the maximum number of active interfaces in a vertex \cite{EnergyConsumptionBalancingInMultiInterfaceNetworks,MinMaxCoverageInMultiInterfaceNetworksPathwidth,OnComputationalComplexityOfCoveringMultiInterfaceNetworks}, yielding efficient algorithms in some graph classes. Some papers also study other problems in multi-interface networks, like cheapest path \cite{ExploitingMultiInterfaceNetworksConnectivityAndCheapestPaths}, matchings \cite{MaximumMatchingInMultiInterfaceNetworks}, or flow problems \cite{FlowProblemsInMultiInterfaceNetworks}.
Closely related combinatorial problems are \textsc{Set Cover}, \textsc{Minimum Spanning Tree} or \textsc{Steiner Tree Forest}, which are used for hardness results \cite{CostMinimizationInWirelessNetworks,MinimizeTheMaximumDutyInMultiInterfaceNetworks,OnComputationalComplexityOfCoveringMultiInterfaceNetworks}, and algorithms \cite{MinimizeTheMaximumDutyInMultiInterfaceNetworks,ExploitingMultiInterfaceNetworksConnectivityAndCheapestPaths,EnergyEfficientCommunicationInMultiInterfaceWirelessNetworks}.

\subsection{Our contribution and techniques}

We study the approximability of the \coverage and \connectivity problems minimizing the max cost under heterogeneous cost functions. We give a set-cover-style LP relaxation for the \coverage problem and show that treating the resulting fractional solution as a probability distribution and applying a randomized rounding procedure yields an $O\br{\log n}$-approximation with high probability. This is the first logarithmic approximation for this problem with heterogeneous costs, and it matches the previously known lower bound of $\Omega\br{\log n}$ for stars with unit costs \cite{MinimizeTheMaximumDutyInMultiInterfaceNetworks}. It also improves the previously best known $O\br{\ln \Delta + b \cdot \min\brc{\ln \Delta, c_{\max}}}$-approximation ratio \cite{MinimizeTheMaximumDutyInMultiInterfaceNetworks} (for some graph input parameter $b$), which was limited to the homogeneous case. 
A simpler deterministic rounding of the same LP gives a $k$-approximation, generalizing the $1+(k-2) \frac{c_{\max}}{2 c_{\min}}$-approximation for the min-max variant with homogeneous costs\cite{MinimizeTheMaximumDutyInMultiInterfaceNetworks}.

For the \connectivity problem, we extend the LP by adding flow-based cut constraints encoding global connectivity. Our rounding procedure is similar to the one for the \coverage problem; however, this time the process is iterative and the algorithm proceeds until all vertices are connected. Each iteration is randomized and we show that in expectation, in each such round the algorithm covers a significant fraction of edges of a graph that spans all vertices. This yields an $O\br{\log^2 n}$-approximation algorithm, which is the first non-trivial approximation algorithm for this problem.

\subsection{Paper Organization}

The rest of the paper is structured as follows:
In \Cref{sec:preliminaries}, we give notations, define the problems formally, and recall the elements of probability theory useful for the analysis of our algorithms. We also describe the preprocessing required for our randomized algorithms for the \coverage and \connectivity problems to work.
In \Cref{sec:coverage}, we present the $O\br{\log n}$-approximation and the $k$-approximation algorithms for the \coverage problem.
In \Cref{sec:connectivity}, we present the $O\br{\log^2 n}$-approximation algorithm for the \connectivity problem.
% We conclude with a short discussion on open questions and future work.